
\section{博特--莫尔斯情形}
\label{sec:8-8}

\liftIndex{I305}本节讨论一般博特--莫尔斯情形下连接轨道空间的定向问题。
首先回顾第~\ref{subsec:3-7-5} 小节的一些记号。设 $L^{(0)}$ 与
$L^{(1)}$ 是两个干净相交的拉格朗日子流形，并以 $R_h$ 表示
$L^{(0)}\cap L^{(1)}$ 的一个连通分支。定义
\[
 V_h=
 \left.
 \frac{TL^{(0)}+TL^{(1)}}
 {(TL^{(0)}+TL^{(1)})^{\perp\omega}}
 \right|_{R_h}.
\]
$V_h$ 是 $R_h$ 上的向量丛；我们把它看作 $TM|_{R_h}$ 的子丛。

对每个 $p\in R_h$，以
\[
 \mathcal P_{R_h}(T_pL^{(0)},T_pL^{(1)})
\]
表示 $T_pM$ 的定向拉格朗日格拉斯曼流形中所有如下道路的空间：
道路形如
\[
 t\longmapsto\lambda(t)\oplus T_pR_h,
\]
并满足
\[
 \lambda(0)\oplus T_pR_h=T_pL^{(0)},\qquad
 \lambda(1)\oplus T_pR_h=T_pL^{(1)}.
\]
这里 $\lambda(t)$ 是 $V_h|_p$ 中拉格朗日子空间的道路。第 3 章考虑的是
\[
 \lambda(t)\oplus T_pR_h=\lambda_{w,\lambda_0}
\]
这一情形。应当指出，$R_h$ 本身未必可定向。记
\[
 \mathcal P_{R_h}(TL^{(0)},TL^{(1)})
 =
 \bigcup_{p\in R_h}
 \mathcal P_{R_h}(T_pL^{(0)},T_pL^{(1)}).
\]

以 $\widetilde\lambda_p$ 表示 $[0,1]\times T_pM$ 中同道路
$\lambda\oplus T_pR_h$ 对应的拉格朗日子丛。

本节所谓连接轨道，是满足下列条件的全纯映射
$u\colon D^2\to M$：
\[
\begin{aligned}
 u(\partial_0D^2)&\subset L^{(0)},&
 u(\partial_1D^2)&\subset L^{(1)},\\
 u(1)&=q\in R_{h'},&
 u(-1)&=p\in R_h.
\end{aligned}
\]
这里 $\partial_0D^2$ 是虚部为负的圆弧，
$\partial_1D^2$ 是虚部为正的圆弧。以
$\mathcal M(R_{h'},R_h)$ 表示这些映射 $u$ 的模空间。
注意，这一记号中 $R_{h'},R_h$ 的次序同第 3.7 节的约定相反
（见第~\ref{subsec:8-1-3} 小节开头）。也就是说，
\[
 \mathcal M(R_{h'},R_h)
 =
 \bigcup_{w'}
 \mathcal M_{0,0}([h,w],[h',w']),
\]
右端采用命题~\ref{prop:3-7-55} 的记号；其中固定 $w$，并对 $w'$ 取并。

回忆
\[
\begin{aligned}
 Z_-&=(D^2\cap\{\operatorname{Re}z\leq0\})
      \cup([0,\infty)\times[0,1]),\\
 Z_+&= ((-\infty,0]\times[0,1])
      \cup(D^2\cap\{\operatorname{Re}z\geq0\}).
\end{aligned}
\]
第~\ref{subsec:3-7-5} 小节定义了多尔博算子
\[
\begin{aligned}
 \bar\partial_{\lambda\oplus TR_h,Z_\pm}\colon&
 W_{\lambda\oplus TR_h}^{1,p}(Z_\pm;T_pM)\\
 &\longrightarrow
 L^p(Z_\pm;T_pM\otimes\Lambda^{0,1}(Z_\pm)).
\end{aligned}
\]
（见定义~\ref{def:3-7-62} 之前。）%
\translatornote{源书把定义域的靶切空间印成 $T_0M$，而同式右端、
道路的基点及第 3.7.5 小节的原定义均为 $T_pM$，故恢复如上。}
算子族
\[
 \left\{
 \bar\partial_{\lambda\oplus TR_h,Z_\pm}
 \right\}_{\lambda\oplus TR_h\in
 \mathcal P_{R_h}(TL^{(0)},TL^{(1)})}
\]
由 $\mathcal P_{R_h}(TL^{(0)},TL^{(1)})$ 参数化。我们将证明，存在纤维丛
\[
 \widetilde{\mathcal I}(R_h)
 \longrightarrow
 \mathcal P_{R_h}(TL^{(0)},TL^{(1)}),
\]
使行列式线丛拉回后下降到 $R_h$。

注~\ref{rem:8-1-15}(1) 在当前情形中推广如下。沿用该处记号，并置
$\lambda_p(t)=\lambda(t)\oplus T_pR_h$，有同构
\[
\begin{aligned}
 \operatorname{Index}
 (\bar\partial_{u_p;\lambda_p,\lambda_p})
 &\cong
 \operatorname{Index}(\bar\partial_{\lambda_p,Z_+})
 \oplus T_pR_h
 \oplus
 \operatorname{Index}(\bar\partial_{\lambda_p,Z_-}),\\
 \operatorname{Index}
 (\bar\partial_{u_p;\lambda_p,\lambda_p})
 &\cong T_pL^{(0)}.
\end{aligned}
\]
因此，
$T_pR_h\oplus
\operatorname{Index}(\bar\partial_{\lambda\oplus TR_h,Z_-})$
的定向确定
$\operatorname{Index}(\bar\partial_{\lambda\oplus TR_h,Z_+})$
的定向。

选定
\[
 \lambda_p\in
 \mathcal P_{R_h}(T_pL^{(0)},T_pL^{(1)}).
\]
把算子 $\bar\partial_{\lambda_p,Z_+}$ 同由
\[
 \lambda'_p\in
 \mathcal P_{R_h}(T_pL^{(0)},T_pL^{(1)})
\]
参数化的算子族 $\bar\partial_{\lambda'_p,Z_-}$ 胶合，得到一族定义在
$D^2$ 上、带全实边界条件的多尔博算子。%
\translatornote{源书在上述两处道路空间的第二个端点都漏排了 $L$，
写成 $T_p^{(1)}$；道路空间的定义及本页此前所有同类记号均要求
$T_pL^{(1)}$。}
同注~\ref{rem:8-1-15}(3) 一样，可以看出，这一算子族的行列式线丛在
$\mathcal P_{R_h}(T_pL^{(0)},T_pL^{(1)})$ 的每个连通分支上都是非平凡的。
由于 $T_pL^{(0)}$ 已定向，而算子
$\bar\partial_{\lambda_p,Z_+}$ 不依赖于 $\lambda'_p$，算子族
$\bar\partial_{\lambda'_p,Z_-}$ 的行列式线丛也是非平凡的。

现在引入第~\ref{subsec:8-1-3} 小节中
$\mathcal I(p)$ 与 $\widetilde{\mathcal I}(p)$ 的族版本。
以 $\mathcal I(R_h)$ 表示道路 $\lambda_p$ 与平凡化
\[
 \sigma\colon[0,1]\times\mathbb R^n
 \longrightarrow\widetilde\lambda_p
\]
所成的空间。对 $\sigma$，考虑覆盖 $\sigma_{t=i*}$ 的提升
\[
 \iota_i\colon
 \frac{\operatorname{Spin}(n)\times\operatorname{Spin}(V_p)}
 {\{\pm1\}}
 \longrightarrow
 P_{\mathrm{spin}}(L^{(i)},V)|_p,
 \qquad i=0,1.
\]
定义 $\widetilde{\mathcal I}(R_h)$ 为四元组
$(\lambda_p,\sigma,\iota_0,\iota_1)$ 所成的空间。存在一列自然投影
\[
 \Pi\colon
 \widetilde{\mathcal I}(R_h)
 \longrightarrow\mathcal I(R_h)
 \longrightarrow\mathcal P_{R_h}(TL^{(0)},TL^{(1)})
 \longrightarrow R_h.
\]

以 $\mathcal D(R_h)^-$ 表示算子族
$\{\bar\partial_{\lambda'_p,Z_-}\}$ 向
$\widetilde{\mathcal I}(R_h)$ 的拉回。于是有下列结果。

\begin{proposition}{8.8.1}\label{prop:8-8-1}
算子族 $\mathcal D(R_h)^-$ 的指标丛之行列式线丛，下降为
$R_h$ 上的一个实线丛。
\end{proposition}

\begin{proof}
注意，
\[
 \Pi^{-1}(p)\subset\widetilde{\mathcal I}(R_h)
\]
不是连通的；这是因为
\[
 \mathcal P_{R_h}(T_pL^{(0)},T_pL^{(1)})
 \subset
 \mathcal P_{R_h}(TL^{(0)},TL^{(1)}),
 \qquad p\in R_h,
\]
不是连通的。因此，先说明如何比较不同连通分支中 $\lambda$ 与
$\lambda'$ 所对应的指标行列式线。

设
\[
 \lambda\oplus T_pR_h
 \in\mathcal P_{R_h}(T_pL^{(0)},T_pL^{(1)}).
\]
如下把秩为 $n$ 的复向量丛 $E\to\mathbb{CP}^1$ 胶合到平凡丛对
$Z_-\times T_pM$ 上，以作扭曲。选取 $T_pM$ 到南极点
$S\in\mathbb{CP}^1$ 上纤维 $E_S$ 的同构，再把
$Z_-\times T_pM$ 在 $O\in Z_-$ 上的纤维同 $E_S$ 识别，便得到一点并
$Z_-\vee\mathbb{CP}^1$ 上的向量丛
\[
 (Z_-\times T_pM)\vee E.
\]
以
\[
 \operatorname{cont}\colon
 Z_-\longrightarrow Z_-\vee\mathbb{CP}^1
\]
表示压缩 $O\in Z_-$ 周围一个小圆周所得的映射，并记
\[
\begin{aligned}
 &(Z_-\times T_pM,\lambda\oplus T_pR_h)
 \#(\mathbb{CP}^1,E)\\
 &\qquad:=
 \operatorname{cont}^*
 \bigl((Z_-\times T_pM,\lambda\oplus T_pR_h)\vee E\bigr).
\end{aligned}
\]

对任意
\[
 \lambda\oplus T_pR_h,\ 
 \lambda'\oplus T_pR_h
 \in\mathcal P_{R_h}(T_pL^{(0)},T_pL^{(1)}),
\]
存在 $\mathbb{CP}^1$ 上的复向量丛 $E$，使
\[
\begin{aligned}
 &(Z_-\times T_pM,\lambda'\oplus T_pR_h)\\
 &\qquad\cong
 (Z_-\times T_pM,\lambda\oplus T_pR_h)
 \#(\mathbb{CP}^1,E).
\end{aligned}
\]
因此，由指标求和公式，
\[
 \operatorname{Index}
 (\bar\partial_{\lambda'\oplus TR_h,Z_-})
 \cong
 \operatorname{Index}(\bar\partial_E)
 \oplus
 \operatorname{Index}
 (\bar\partial_{\lambda\oplus TR_h,Z_-}).
\]
作为复虚向量丛，多尔博算子 $\bar\partial_E$ 的指标带有典范定向。
所以
$\operatorname{Index}
 (\bar\partial_{\lambda'\oplus TR_h,Z_-})$
的定向以不依赖于 $E$ 的方式确定
$\operatorname{Index}
 (\bar\partial_{\lambda\oplus TR_h,Z_-})$
的定向。

下面证明，$\mathcal D(R_h)^-$ 的行列式线丛限制到
$\Pi^{-1}(p)$ 的每个连通分支上都是平凡的。结合上述论证，该行列式线丛
便下降到 $R_h$。%
\translatornote{源书本句省去了此前始终保留的上标 $\,{}^{-}$，写成
$\mathcal D(R_h)$；依命题所定义的算子族补回。}
（由于行列式线丛是实线丛，其结构群可约化为 $\{\pm1\}$。）
选定
\[
 (\lambda_p,\sigma,\iota_0,\iota_1)\in\Pi^{-1}(p).
\]
考虑经
$\mathcal P_{R_h}(T_pL^{(0)},T_pL^{(1)})$
由
\[
 (\lambda'_p,\sigma',\iota'_0,\iota'_1)
 \in\Pi^{-1}(p)
\]
参数化的算子族
\[
 \mathcal D_p=\{\bar\partial_{\lambda'_p,Z_-}\}.
\]
把 $\mathcal D_p$ 同 $\bar\partial_{\lambda_p,Z_+}$ 胶合，得到
$D^2$ 上的一族多尔博算子。由于
$(\lambda_p,\sigma,\iota_0,\iota_1)$ 与
$(\lambda'_p,\sigma',\iota'_0,\iota'_1)$ 以相容方式确定全实子丛族上的
spin 结构，这族多尔博算子的行列式线丛在 $\Pi^{-1}(p)$ 上是平凡的。
因此 $\mathcal D_p$ 的行列式线丛也是平凡的。命题得证。
\end{proof}

\begin{definition}{8.8.2}\label{def:8-8-2}
以 $\Theta^\pm_{R_h}$ 表示命题~\ref{prop:8-8-1} 从算子
$\bar\partial_{\lambda\oplus TR_h,Z_\pm}$ 的指标丛行列式所得的
$R_h$ 上局部系统。
\end{definition}

在有限维博特--莫尔斯理论中，局部系统 $\Theta^+_{R_h}$ 与
$\Theta^-_{R_h}$ 分别对应于黑塞算子在法丛上的限制之正定部分
（稳定方向）与负定部分（不稳定方向）的定向丛。

注意，$R_h$ 未必可定向，连接轨道空间也可能不可定向。因此必须处理以
不可定向空间为底的纤维积。为此，下面引入带仓西结构空间的定向丛，并
给出它的基本性质。\liftIndex{I306}

\begin{definition}{8.8.3}\label{def:8-8-3}
对带切丛仓西结构的局部坐标
$(s;E\to U)$，称
\[
 \det E^*\otimes\det TU
\]
为其\emph{定向丛}。
\end{definition}

若仓西结构带有切丛，则局部坐标（仓西邻域）的定向丛自然地粘合成带仓西
结构空间上的一个线丛（见定义 A1.14）。称它为该带仓西结构空间的定向丛。

下面规定纤维积定向丛的识别方式。

\begin{statement}{约定}{8.8.4}\label{conv:8-8-4}
设 $O_X,O_Y$ 分别是带切丛仓西结构空间 $X,Y$ 的定向丛。设
\[
 f\colon X\to B,\qquad g\colon Y\to B
\]
是到流形 $B$ 的弱浸没强连续映射，并以 $O_B$ 表示 $B$ 的定向丛。

把纤维积 $X_f\times_gY$ 的定向丛
$O_{X_f\times_gY}$ 识别为
\[
 \operatorname{pr}_X^*(O_X\otimes f^*O_B)
 \otimes\pi^*O_B
 \otimes
 \operatorname{pr}_Y^*(g^*O_B\otimes O_Y).
\]
上述实线丛也自然地同
\[
 \operatorname{pr}_X^*O_X
 \otimes\pi^*O_B
 \otimes\operatorname{pr}_Y^*O_Y
\]
识别。这里 $\operatorname{pr}_X,\operatorname{pr}_Y$ 分别是到
$X,Y$ 的投影，而 $\pi$ 是到 $B$ 的投影。

必须强调，上述张量积中因子的次序至关重要。例如，当 $X,Y$ 是流形时，
上述识别对应于同构
\[
 O_X\otimes O_Y\longrightarrow O_{X\times Y},
 \qquad s\otimes t\longmapsto s\wedge t.
\]
\end{statement}

像引理~\ref{lem:8-2-3} 一样，可以证明下列性质。\translatornote{源书在下列
(2) 的左端把最后一个因子重复写成 $X_2$，而右端及结合律的型别均要求它为
$X_3$。}

\begin{lemma}{8.8.5}\label{lem:8-8-5}
有下列识别。
\begin{enumerate}[label=(\arabic*),leftmargin=2.8em]
\item
\[
\begin{aligned}
 O_{\partial(X\times_BY)}|_{\partial X\times_BY}
 &\cong O_{\partial X\times_BY},\\
 O_{\partial(X\times_BY)}|_{X\times_B\partial Y}
 &\cong(-1)^{x+y}O_{X\times_B\partial Y}.
\end{aligned}
\]

\item
\[
 O_{(X_1\times_{B_1}X_2)\times_{B_2}X_3}
 \cong
 O_{X_1\times_{B_1}(X_2\times_{B_2}X_3)}.
\]

\item
\[
\begin{aligned}
 O_{X_1\times_{B_1\times B_2}(X_2\times X_3)}
 \cong
 (-1)^{b_2(b_1+x_2)}
 O_{(X_1\times_{B_1}X_2)\times_{B_2}X_3}.
\end{aligned}
\]
更一般地，
\[
\begin{aligned}
 &O_{X_1\times_{B_1\times\cdots\times B_\ell}
 (X_2\times\cdots\times X_{\ell+1})}\\
 &\quad\cong
 (-1)^{
 \sum_{k=2}^{\ell}b_k
 \sum_{j=2}^{k}(b_{j-1}+x_j)}
 O_{(\cdots(X_1\times_{B_1}X_2)\cdots)
 \times_{B_\ell}X_{\ell+1}}.
\end{aligned}
\]
\end{enumerate}
\end{lemma}

下面用引理~\ref{lem:8-8-5} 研究本节开头所引入的模空间
$\mathcal M(R_{h'},R_h)$ 的定向丛，并解决引理~\ref{lem:3-7-70}
中的定向问题。

\begin{proposition}{8.8.6}\label{prop:8-8-6}
拉格朗日子流形对 $(L^{(0)},L^{(1)})$ 的相对 spin 结构典范地诱导同构
\[
 O_{\widetilde{\mathcal M}(R_{h'},R_h)}
 \cong
 \operatorname{ev}_0^*\Theta^+_{R_{h'}}
 \otimes
 \operatorname{ev}_1^*\Theta^-_{R_h}.
\]
此外，还有典范同构
\[
 O_{\widetilde{\mathcal M}(R_{h'},R_h)}
 \cong
 \pi^*O_{\mathcal M(R_{h'},R_h)}
 \otimes O_{\mathbb R},
\]
其中
\[
 \pi\colon
 \widetilde{\mathcal M}(R_{h'},R_h)
 \longrightarrow
 \mathcal M(R_{h'},R_h)
\]
是自然投影，而 $O_{\mathbb R}$ 是平移 $\mathbb R$-作用的轨道定向。%
\translatornote{源书把自然投影本身记作 $\pi^*$；但上一式已经用
$\pi^*$ 表示拉回，故投影应为 $\pi$。}
因而
\[
 O_{\mathcal M(R_{h'},R_h)}
 \cong
 \operatorname{ev}_0^*\Theta^+_{R_{h'}}
 \otimes
 \operatorname{ev}_1^*\Theta^-_{R_h}
 \otimes O_{\mathbb R}.
\]
\end{proposition}

\begin{proof}
可以假定评价映射
\[
 \operatorname{ev}_0,\operatorname{ev}_1\colon
 \mathcal M(R_{h'},R_h)\longrightarrow L
\]
是弱浸没的；这由命题~\ref{prop:7-1-1} 的博特--莫尔斯版本得到。

取 $u\in\mathcal M(R_{h'},R_h)$，以
$D_u\bar\partial$ 表示柯西--黎曼算子的线性化。（当
$R_h=R_{h'}=L$ 时，定义 $D_u\bar\partial$ 所需的弗雷德霍姆理论见
第 7.1.2 小节的引理~\ref{lem:7-1-10}；把它推广到一般博特--莫尔斯情形
是直接的。）通过把 $T_pR_h$（分别地，$T_qR_{h'}$）同
$J\frac d{dt}$ 在 $z_1$（分别地，$z_0$）处的零特征空间识别，得到从
$D_u\bar\partial$ 的定义域到
$T_pR_h\oplus T_qR_{h'}$ 的同态。

令
\[
 p'=u(+\infty,t)\in R_{h'},\qquad
 p=u(-\infty,t)\in R_h,
\]
并取
\[
\begin{aligned}
 \lambda_p\oplus T_pR_h
 &\in\mathcal P_{R_h}(T_pL^{(0)},T_pL^{(1)}),\\
 \lambda_{p'}\oplus T_{p'}R_{h'}
 &\in\mathcal P_{R_{h'}}(T_{p'}L^{(0)},T_{p'}L^{(1)}).
\end{aligned}
\]
把 $u^*TM$ 平凡化，再通过以显然方式胶合
\[
 \tau\longmapsto T_{u(\tau,i)}L^{(i)},\qquad
 \lambda_p\oplus T_pR_h,\qquad
 \lambda_{p'}\oplus T_{p'}R_{h'}
\]
定义这个平凡丛纤维中拉格朗日线性子空间的道路
$\lambda(\lambda_{p'},u,\lambda_p)$。%
\translatornote{源书在上述两处道路成员关系及随后的胶合资料中写
$\lambda_p\oplus R_h$、$\lambda_{p'}\oplus R_{h'}$；由于 $R_h$ 是
流形而 $\lambda_p$ 是线性子空间，依本节开头定义补回相应切空间。}

于是，由指标求和公式的族版本，得到下列指标纤维积的同构。（纤维积定向
的约定见约定~\ref{conv:8-2-1}(4)，指标丛纤维积的定义见
式~\eqref{eq:8-3-9}。）
\[
\begin{aligned}
 &\operatorname{Index}
 \bigl(\bar\partial_{\lambda(\lambda_{p'},u,\lambda_p)}\bigr)\\
 &\quad\cong
 \left(
 \operatorname{Index}
 (\bar\partial_{\lambda_{p'}\oplus TR_{h'},Z_+})
 \oplus T_{p'}R_{h'}
 \right)
 \mathop{\times}_{T_{p'}R_{h'}}
 \operatorname{Index}(D_u\bar\partial)\\
 &\hspace{9em}{}
 \mathop{\times}_{T_pR_h}
 \left(
 T_pR_h\oplus
 \operatorname{Index}
 (\bar\partial_{\lambda_p\oplus TR_h,Z_-})
 \right).
\end{aligned}
\]
同定理~\ref{thm:8-1-1}（以及定理~\ref{thm:8-1-14}）的证明一样，可以
证明
$\operatorname{Index}
(\bar\partial_{\lambda(\lambda_{p'},u,\lambda_p)})$
带有典范定向。

因此，$\operatorname{Index}(D_u\bar\partial)$ 的定向由下列各空间的定向
确定：
\[
\begin{gathered}
 \operatorname{Index}
 (\bar\partial_{\lambda_{p'}\oplus T_{p'}R_{h'},Z_+})
 \oplus T_{p'}R_{h'},\qquad
 T_{p'}R_{h'},\qquad T_pR_h,\\
 T_pR_h\oplus
 \operatorname{Index}
 (\bar\partial_{\lambda_p\oplus T_pR_h,Z_-}).
\end{gathered}
\]
由于 $T_pR_h$ 与 $T_{p'}R_{h'}$ 都各出现两次，
$\operatorname{Index}(D_u\bar\partial)$ 的定向不依赖于
$T_{p'}R_{h'}$、$T_pR_h$ 的定向。由此得到命题的前半；后半立即由前半
及约定~\ref{conv:8-2-1}(2) 得到。
\end{proof}

回忆，我们采用上同调约定。现作如下规定：标准单形取通常的定向。设
$(\Delta,\sigma)$ 是未必可定向的流形 $X$ 中的奇异单形，并带有定向丛
$O_X=\det TX$ 中的系数。例如，法丛已定向的嵌入子流形表示一个普通
上链。

下面说明如何在我们的构造中使用带局部系数的链。设 $S$ 是 $R_h$ 中
带
\[
 O_{R_h}\otimes\Theta^-_{R_h}
 =\det TR_h\otimes\Theta^-_{R_h}
 \cong\Theta^+_{R_h}
\]
系数的奇异单形。回忆，连接轨道（即
$\mathcal M(R_{h'},R_h)$ 的元素）是伪全纯映射 $u\colon D^2\to M$，
满足
\[
\begin{gathered}
 u(\partial_0D^2)\subset L^{(0)},\qquad
 u(\partial_1D^2)\subset L^{(1)},\\
 u(1)\in R_{h'},\qquad u(-1)\in R_h,
\end{gathered}
\]
其中 $\partial_0D^2$ 是虚部为负的圆弧，
$\partial_1D^2$ 是虚部为正的圆弧。

\begin{proposition}{8.8.7}\label{prop:8-8-7}
设 $S$ 是带 $O_{R_h}\otimes\Theta^-_{R_h}$ 系数的奇异单形，则
\[
 \bigl(\mathcal M(R_{h'},R_h)\times_{R_h}S,\operatorname{ev}_0\bigr)
\]
是带 $O_{R_{h'}}\otimes\Theta^-_{R_{h'}}$ 系数的链。
\end{proposition}

\begin{proof}
只须考虑 $S$ 是带 $O_{R_h}\otimes\Theta^-_{R_h}$ 中系数的奇异单形
这一情形。写成 $S=(\phi,s)$，其中
\[
 \phi\colon\Delta^k\longrightarrow R_h,
\]
而 $s$ 是
$\phi^*(O_{R_h}\otimes\Theta^-_{R_h})$ 的非零平坦截面。注意，$s$
给出
\[
 T_pR_h\oplus
 \operatorname{Index}
 (\bar\partial_{\lambda_p\oplus TR_h,Z_-})
\]
的定向。现在把 $\mathcal M(R_{h'},R_h)$ 与 $S$ 在 $R_h$ 上取纤维积。
由于 $O_{\Delta^k}$ 已由标准定向平凡化，可把 $S$ 看作配备局部系统
\[
 \phi^*(O_{R_h}\otimes\Theta^-_{R_h})\otimes O_{\Delta^k}
\]
之非零平坦截面的奇异单形。（现阶段不需要最后这个已平凡化的因子
$O_{\Delta^k}$，但它使论证更为典范；见证明末尾。）由约定
\ref{conv:8-8-4} 与命题~\ref{prop:8-8-6}，
\[
\begin{aligned}
 O_{\mathcal M(R_{h'},R_h)\times_{R_h}S}
 &\cong
 \operatorname{ev}_0^*\Theta^+_{R_{h'}}
 \otimes\operatorname{ev}_1^*\Theta^-_{R_h}
 \otimes O_{\mathbb R}
 \otimes\operatorname{ev}_1^*O_{R_h}
 \otimes O_{\Delta^k}\\
 &\cong
 (-1)^{r_h}
 \operatorname{ev}_0^*\Theta^+_{R_{h'}}
 \otimes
 \operatorname{ev}_1^*(\Theta^-_{R_h}\otimes O_{R_h})
 \otimes O_{\mathbb R}\otimes O_{\Delta^k},
\end{aligned}
\]
其中 $r_h=\dim R_h$。所以
\[
\begin{aligned}
 &\operatorname{ev}_0^*\Theta^+_{R_{h'}}
 \otimes
 O_{\mathcal M(R_{h'},R_h)\times_{R_h}S}\\
 &\qquad\cong
 (-1)^{r_h}
 \operatorname{ev}_1^*(\Theta^-_{R_h}\otimes O_{R_h})
 \otimes O_{\mathbb R}\otimes O_{\Delta^k}.
\end{aligned}
\]
注意，$O_{\mathbb R}$ 与 $O_{\Delta^k}$ 都典范地定向。又有
\[
 \operatorname{ev}_1^*(\Theta^-_{R_h}\otimes O_{R_h})
 \cong
 (-1)^{r_h\cdot\mu(h,\lambda_p)}
 \operatorname{ev}_1^*(O_{R_h}\otimes\Theta^-_{R_h}),
\]
其中
\[
 \mu(h,\lambda_p)
 =
 \operatorname{Index}
 (\bar\partial_{\lambda_p\oplus TR_h,Z_-}).
\]
由于 $\operatorname{ev}_1$ 与 $\phi$ 在
$\mathcal M(R_{h'},R_h)\times_{R_h}S$ 上相容，平坦截面 $s$ 诱导
\[
 \operatorname{ev}_0^*\Theta^+_{R_{h'}}
 \otimes O_{\mathcal M(R_{h'},R_h)\times_{R_h}S}
\]
的平坦截面。因此，对该纤维积作扰动并三角剖分后，得到一条带
\[
 O_{R_{h'}}\otimes\Theta^-_{R_{h'}}
 \cong\Theta^+_{R_{h'}}
\]
系数的链。
\end{proof}

计入命题~\ref{prop:8-8-7} 后，从现在起可以不再显式写出局部系数的
影响，而把计算当作使用普通有向链来进行。为了定义滤过
$A_\infty$ 双模运算，我们还要进一步修正定向（见定义 8.8.11）。

下文以 $P_1,\ldots,P_k$ 表示 $L^{(0)}$ 或 $L^{(1)}$ 中的光滑奇异
单形。此时，无穷圆柱 $\mathbb R\times[0,1]$ 边界上标记点的次序任意，
不必服从标准次序。也就是说，不对 $\mathbb R\times\{0,1\}$ 上标记点的
次序作任何特别规定；而且 $P_i\subset L^{(0)}$ 还是
$P_i\subset L^{(1)}$，同次序 $i$ 没有关系。把这个模空间记为
\begin{equation}
\begin{aligned}
 &\mathcal M_{k+2}((R_{h'},w'),(R_h,w))
 \mathop{\times}_{R_h\times L\times\cdots\times L}
 (S\times P_1\times\cdots\times P_k).
\end{aligned}
\tag{*}\label{eq:8-8-star}
\end{equation}
见图~\ref{fig:8-8-1}。注意，当
\[
 (P_1,\ldots,P_k)
 =
 (P_1^{(1)},\ldots,P_{k_1}^{(1)},
  P_1^{(0)},\ldots,P_{k_0}^{(0)})
\]
且 $P_j^{(i)}$ 是 $L^{(i)}$ 中的光滑奇异单形时，模空间
\eqref{eq:8-8-star} 包含如下模空间作为开子集：
\[
\begin{aligned}
 \mathcal M_{k_1,k_0,(+\infty)}
 \bigl(
 L^{(1)},L^{(0)};[h',w'],[h,w];
 \vec P^{\,(1)},S,\vec P^{\,(0)}
 \bigr).
\end{aligned}
\]
它在注~\ref{rem:3-7-58} 之后定义，并在定义~\ref{def:3-7-71} 中用于
定义滤过 $A_\infty$ 双模结构。（再次提醒，$[h,w]$ 与 $[h',w']$ 的
次序在这里是反过来的。）现在要给模空间 \eqref{eq:8-8-star} 赋予定向。

\begin{figure}[htbp]
\centering
\setlength{\unitlength}{1mm}
\begin{picture}(126,55)
  \linethickness{0.48pt}
  \qbezier(17,26)(59,52)(105,26)
  \qbezier(17,26)(59,0)(105,26)
  \put(7,36){\line(1,-1){20}}
  \put(7,16){\line(1,1){20}}
  \put(95,16){\line(1,1){20}}
  \put(95,36){\line(1,-1){20}}
  \put(30,32.8){\makebox(0,0){$\times$}}
  \put(47,37.9){\makebox(0,0){$\times$}}
  \put(74,37.6){\makebox(0,0){$\times$}}
  \put(91,32.7){\makebox(0,0){$\times$}}
  \put(26,39){\makebox(0,0){$P_8$}}
  \put(45,42){\makebox(0,0){$P_3$}}
  \put(76,44){\makebox(0,0){$P_2$}}
  \put(95,38){\makebox(0,0){$P_7$}}
  \put(35,17.3){\makebox(0,0){$\times$}}
  \put(53,13.3){\makebox(0,0){$\times$}}
  \put(70,13.3){\makebox(0,0){$\times$}}
  \put(87,17.6){\makebox(0,0){$\times$}}
  \put(32,10.5){\makebox(0,0){$P_6$}}
  \put(52,7){\makebox(0,0){$P_1$}}
  \put(71,6.5){\makebox(0,0){$P_5$}}
  \put(92,9.5){\makebox(0,0){$P_4$}}
  \put(4,26){\makebox(0,0){$S$}}
  \put(120,37){\makebox(0,0)[l]{$L^{(0)}$}}
  \put(120,15){\makebox(0,0)[l]{$L^{(1)}$}}
\end{picture}
\par\smallskip
\renewcommand{\thefigure}{8.8.1}
\refstepcounter{figure}
\textbf{\figurename~\thefigure}\label{fig:8-8-1}
\end{figure}


设 $S$ 是 $R_h$ 中的奇异单形。定义其次数为
\begin{equation}
 \deg S=\mu([h,w])+\dim R_h-\dim S.
\tag{8.8.8}\label{eq:8-8-8}
\end{equation}
（见定义~\ref{def:3-7-52} 之前。）注意，对任意
\[
 \lambda\oplus T_pR_h
 \in\mathcal P_{R_h}(T_pL^{(0)},T_pL^{(1)}),
\]
都有
\[
 \mu([h,w])
 \equiv
 \operatorname{Index}
 (\bar\partial_{\lambda\oplus TR_h,Z_-})
 \pmod2.
\]
这是由于右端模 $2$ 意义下的独立性，以及
$\mu([h,w])$ 的定义、定义~\ref{def:3-7-62} 与引理
\ref{lem:3-7-69}。%
\translatornote{源书在上述指标式的右端多排了一个不匹配的右括号；
译文按配对规范化。}

为便于计算，先引入如下临时约定；最终的定向约定将在定义 8.8.11 给出。
置
\[
\begin{aligned}
 &\mathcal M_1
 ((R_{h'},w'),(R_h,w);S,P_1,\ldots,P_k)\\
 &\quad=
 (-1)^\epsilon
 \mathcal M_{k+2}((R_{h'},w'),(R_h,w))
 \mathop{\times}_{R_h\times L\times\cdots\times L}
 (S\times P_1\times\cdots\times P_k),
\end{aligned}
\]
其中
\[
\begin{aligned}
 \epsilon={}&
 k(n+1)\mu([h,w])+\mu([h',w'])
 +k(n+1)\deg S\\
 &{}+(n+1)
 \sum_{j=1}^{k-1}\sum_{i=1}^{j}\deg P_i.
\end{aligned}
\]

\begin{remark}{8.8.9}\label{rem:8-8-9}
同时交换标记点 $z_i,z_{i+1}$（其中 $i>1$；注意 $z_1=-\infty$）与
$P_{i-1},P_i$ 时，有同构
\[
\begin{aligned}
 &\mathcal M_1
 ((R_{h'},w'),(R_h,w);
 S,P_1,\ldots,P_{i-1},P_i,\ldots,P_k)\\
 &\quad\longrightarrow
 \mathcal M_1
 ((R_{h'},w'),(R_h,w);
 S,P_1,\ldots,P_i,P_{i-1},\ldots,P_k),
\end{aligned}
\]
它是
\[
 (-1)^{(\deg P_{i-1}+1)(\deg P_i+1)}
\]
-定向同构。%
\translatornote{源书的指数第二因子写成 $\deg P_{i+1}+1$；但所述同构
实际交换的是 $P_{i-1}$ 与 $P_i$，且引理~\ref{lem:8-4-3} 对这两个
相邻输入给出的柯西符号正是上式，故改正。}
证明同引理~\ref{lem:8-4-3} 的证明类似。
\end{remark}

从现在起省略 $w,w'$ 等资料，简写为
\[
 \mathcal M_1(R_{h'},R_h;S,P_1,\ldots,P_k)
 =
 \mathcal M_1
 ((R_{h'},w'),(R_h,w);S,P_1,\ldots,P_k).
\]
也以 $\mu(R_h)$ 表示 $\mu([h,w])\bmod2$。

$\mathcal M_1(R_{h'},R_h;S,P_1,\ldots,P_k)$ 有四类边界分支。第一类
出现在把 $S$ 换成 $\partial S$ 时；第二类出现在把 $P_i$ 换成
$\partial P_i$ 时；第三类出现在连接轨道分裂成断裂连接轨道时。交换
标记点与各 $P_i$ 后，第三类写成
\[
\begin{aligned}
 \mathcal M_1\bigl(
 R_{h'},R_{h''};
 \mathcal M_1(R_{h''},R_h;S,P_1,\ldots,P_{k_1}),
 P_{k_1+1},\ldots,P_{k_1+k_2}
 \bigr).
\end{aligned}
\]
在下面定理的 (3) 中，置 $k=k_1$、$\ell=k_2$、
$Q_i=P_{k_1+i}$，把上述模空间写成
\[
\begin{aligned}
 \mathcal M_1\bigl(
 R_{h'},R_{h''};
 \mathcal M_1(R_{h''},R_h;S,P_1,\ldots,P_k),
 Q_1,\ldots,Q_\ell
 \bigr).
\end{aligned}
\]
第四类出现在全纯圆盘冒泡时。

\begin{theorem}{8.8.10}\label{thm:8-8-10}
设 $R_{h_1},R_{h_2},R_{h_3}$ 是
$L^{(0)}\cap L^{(1)}$ 的连通分支。%
\translatornote{源书此处误写为“connecting components”；本节开头
定义的对象及下标 $h_i$ 都是交集的 connected components。}
记
\[
 r_i=\dim R_{h_i},\qquad
 \mu_i\equiv\mu(R_{h_i})\pmod2.
\]
则有下列关系。
\begin{enumerate}[label=(\arabic*),leftmargin=2.8em]
\item
\[
\begin{aligned}
 &(-1)^{\epsilon_1}
 \mathcal M_1
 (R_{h_2},R_{h_1};\partial S,P_1,\ldots,P_k)\\
 &\qquad\subset
 \partial\mathcal M_1
 (R_{h_2},R_{h_1};S,P_1,\ldots,P_k),
\end{aligned}
\]
其中
\[
 \epsilon_1=(\mu_2+r_2)+(\mu_1+r_1)+1.
\]

\item
\[
\begin{aligned}
 &(-1)^{\epsilon_2}
 \mathcal M_1
 (R_{h_2},R_{h_1};
 S,P_1,\ldots,\partial P_i,\ldots,P_k)\\
 &\qquad\subset
 \partial\mathcal M_1
 (R_{h_2},R_{h_1};S,P_1,\ldots,P_i,\ldots,P_k),
\end{aligned}
\]
其中
\[
 \epsilon_2=(\mu_2+r_2)+n+(\deg S+1)
 +\sum_{j=1}^{i-1}(\deg P_j+1)+1.
\]

\item
\[
\begin{aligned}
 &(-1)^{\epsilon_3}
 \mathcal M_1\bigl(
 R_{h_3},R_{h_2};
 \mathcal M_1(R_{h_2},R_{h_1};S,P_1,\ldots,P_k),
 Q_1,\ldots,Q_\ell
 \bigr)\\
 &\qquad\subset
 \partial\mathcal M_1
 (R_{h_3},R_{h_1};
 S,P_1,\ldots,P_k,Q_1,\ldots,Q_\ell),
\end{aligned}
\]
其中
\[
 \epsilon_3=(\mu_3+r_3)+1.
\]

\item
\[
\begin{aligned}
 &(-1)^{\epsilon_4}
 \mathcal M_1\bigl(
 R_{h_2},R_{h_1};
 S,\mathcal M_{k+1}(\beta;P_1,\ldots,P_k),
 Q_1,\ldots,Q_\ell
 \bigr)\\
 &\qquad\subset
 \partial\mathcal M_1
 (R_{h_2},R_{h_1};
 S,P_1,\ldots,P_k,Q_1,\ldots,Q_\ell),
\end{aligned}
\]
其中
\[
 \epsilon_4=(\mu_2+r_2)+\deg S.
\]
\end{enumerate}
\end{theorem}

\begin{proof}
首先注意
\[
 \dim\mathcal M(R_{h_2},R_{h_1})
 =\mu_2-\mu_1+r_2-1.
\]

\textup{(1)}
由 \eqref{eq:8-8-8}，
\[
\begin{aligned}
 &\mathcal M_1
 (R_{h_2},R_{h_1};\partial S,P_1,\ldots,P_k)\\
 &\quad=
 (-1)^{a_1}
 \mathcal M_{k+2}(R_{h_2},R_{h_1})
 \mathop{\times}_{R_{h_1}\times L\times\cdots\times L}
 \left(\partial S\times\prod_{i=1}^{k}P_i\right),
\end{aligned}
\]
其中
\[
 a_1=
 k\mu_1(n+1)+\mu_2+k(n+1)\deg\partial S
 +(n+1)\sum_{j=1}^{k-1}\sum_{i=1}^{j}\deg P_i.
\]
此外，
\[
\begin{aligned}
 &\mathcal M_1
 (R_{h_2},R_{h_1};S,P_1,\ldots,P_k)\\
 &\quad=
 (-1)^{a_2}
 \mathcal M_{k+2}(R_{h_2},R_{h_1})
 \mathop{\times}_{R_{h_1}\times L\times\cdots\times L}
 \left(S\times\prod_{i=1}^{k}P_i\right),
\end{aligned}
\]
其中
\[
 a_2=
 k\mu_1(n+1)+\mu_2+k(n+1)\deg S
 +(n+1)\sum_{j=1}^{k-1}\sum_{i=1}^{j}\deg P_i.
\]
由引理~\ref{lem:8-2-3}(1)，
\[
\begin{aligned}
 &(-1)^{a_3}
 \mathcal M_{k+2}(R_{h_2},R_{h_1})
 \mathop{\times}_{R_{h_1}\times L\times\cdots\times L}
 \left(\partial S\times\prod_{i=1}^{k}P_i\right)\\
 &\quad\subset
 \partial\left[
 \mathcal M_{k+2}(R_{h_2},R_{h_1})
 \mathop{\times}_{R_{h_1}\times L\times\cdots\times L}
 \left(S\times\prod_{i=1}^{k}P_i\right)
 \right],
\end{aligned}
\]
其中
\[
 a_3=(r_2+\mu_2)-\mu_1+k-1+r_1+kn.
\]
因此
\[
 \epsilon_1\equiv a_1+a_2+a_3
 \equiv(r_2+\mu_2)+(r_1+\mu_1)+1.
\]

\textup{(2)}
类似地，
\[
\begin{aligned}
 &\mathcal M_1
 (R_{h_2},R_{h_1};
 S,P_1,\ldots,\partial P_i,\ldots,P_k)\\
 &\quad=
 (-1)^{b_1}
 \mathcal M_{k+2}(R_{h_2},R_{h_1})
 \mathop{\times}_{R_{h_1}\times L\times\cdots\times L}
 \left(S\times\prod_{j=1}^{k}P'_j\right),
\end{aligned}
\]
其中
\[
 P'_j=P_j\quad(j\ne i),\qquad P'_i=\partial P_i,
\]
并且
\[
 b_1=
 k\mu_1(n+1)+\mu_2+k(n+1)\deg S
 +(n+1)\sum_{j=1}^{k-1}\sum_{h=1}^{j}\deg P'_h.
\]
注意
\[
 (-1)^cS\times\prod_{j=1}^{k}P'_j
 \subset
 \partial\left(S\times\prod_{j=1}^{k}P_j\right),
\qquad
 c=\dim S+\sum_{j=1}^{i-1}\dim P_j.
\]
由引理~\ref{lem:8-2-3}(1)，
\[
\begin{aligned}
 &(-1)^{b_3}
 \mathcal M_{k+2}(R_{h_2},R_{h_1})
 \mathop{\times}_{R_{h_1}\times L\times\cdots\times L}
 \left(S\times\prod_{j=1}^{k}P'_j\right)\\
 &\quad\subset
 \partial\left[
 \mathcal M_{k+2}(R_{h_2},R_{h_1})
 \mathop{\times}_{R_{h_1}\times L\times\cdots\times L}
 \left(S\times\prod_{j=1}^{k}P_j\right)
 \right],
\end{aligned}
\]
其中
\[
 b_3=a_3+c
 =(r_2+\mu_2)-\mu_1+k-1+r_1+kn
 +\dim S+\sum_{j=1}^{i-1}\dim P_j.
\]
所以
\[
 \epsilon_2\equiv b_1+a_2+b_3
 \equiv
 (r_2+\mu_2)+n+(\deg S+1)
 +\sum_{j=1}^{i-1}(\deg P_j+1)+1.
\]

\textup{(3)}
依连接轨道线性化算子指标的定向约定以及约定
\ref{conv:8-2-1}(4)，可知
\[
 \widetilde{\mathcal M}(R_{h_3},R_{h_1})
 =
 \widetilde{\mathcal M}(R_{h_3},R_{h_2})
 \mathop{\times}_{R_{h_2}}
 \widetilde{\mathcal M}(R_{h_2},R_{h_1}).
\]
于是，命题~\ref{prop:8-3-3} 的类似结果成立。也就是说，像第 8.3 节
那样把 $\ell$ 个标记点置于 $k$ 个标记点之前，便有
\[
\begin{aligned}
 &(-1)^{d_0}
 \mathcal M_{k+2}(R_{h_3},R_{h_2})
 \mathop{\times}_{R_{h_2}}
 \mathcal M_{\ell+2}(R_{h_2},R_{h_1})\\
 &\qquad\subset
 \partial\mathcal M_{k+\ell+2}(R_{h_3},R_{h_1}),
\end{aligned}
\]
其中
\[
 d_0=k\ell+k(\mu_2-\mu_1)+(\mu_3-\mu_2)+k-1+r_3.
\]

依定义，
\[
\begin{aligned}
 &\mathcal M_1\bigl(
 R_{h_3},R_{h_2};
 \mathcal M_1(R_{h_2},R_{h_1};S,P_1,\ldots,P_\ell),
 Q_1,\ldots,Q_k
 \bigr)\\
 &\quad=
 (-1)^{d_1}
 \Bigl[
 \mathcal M_{k+2}(R_{h_3},R_{h_2})\\
 &\hspace{5em}{}
 \mathop{\times}_{R_{h_2}\times L\times\cdots\times L}
 \left(
 \mathcal M_1(R_{h_2},R_{h_1};S,P_1,\ldots,P_\ell)
 \times\prod_{i=1}^{k}Q_i
 \right)
 \Bigr],
\end{aligned}
\]
其中
\[
\begin{aligned}
 d_1={}&
 k(n+1)\mu_2+\mu_3\\
 &{}+k(n+1)
 \left\{
 \deg\mathcal M_1
 (R_{h_2},R_{h_1};S,P_1,\ldots,P_\ell)
 \right\}\\
 &{}+(n+1)
 \sum_{i=1}^{k-1}\sum_{j=1}^{i}\deg Q_j.
\end{aligned}
\]
注意
\[
\begin{aligned}
 &\deg\mathcal M_1
 (R_{h_2},R_{h_1};S,P_1,\ldots,P_\ell)\\
 &\quad=
 r_2+\mu_2-
 \left\{
 r_2+\mu_2-\mu_1+\ell-1
 -(\deg S-\mu_1)-\sum_i\deg P_i
 \right\}\\
 &\quad=
 -\ell+1+\deg S+\sum_i\deg P_i.
\end{aligned}
\]
再由定义，
\[
\begin{aligned}
 &\mathcal M_1(R_{h_2},R_{h_1};S,P_1,\ldots,P_\ell)\\
 &\quad=
 (-1)^{d_2}
 \mathcal M_{\ell+2}(R_{h_2},R_{h_1})
 \mathop{\times}_{R_{h_1}\times L\times\cdots\times L}
 \left(S\times\prod_{i=1}^{\ell}P_i\right),
\end{aligned}
\]
其中
\[
 d_2=
 \ell(n+1)\mu_1+\mu_2+\ell(n+1)\deg S
 +(n+1)\sum_{j=1}^{\ell-1}\sum_{i=1}^{j}\deg P_i.
\]

反复应用引理~\ref{lem:8-2-3}(3)，得到
\[
\begin{aligned}
 &\mathcal M_{k+2}(R_{h_3},R_{h_2})\\[-2pt]
 &\quad
 \mathop{\times}_{R_{h_2}\times L\times\cdots\times L}
 \Biggl[
 \Bigl(
 \mathcal M_{\ell+2}(R_{h_2},R_{h_1})\\
 &\hspace{5em}{}
 \mathop{\times}_{R_{h_1}\times L\times\cdots\times L}
 \left(S\times\prod_{i=1}^{\ell}P_i\right)
 \Bigr)
 \times\prod_{j=1}^{k}Q_j
 \Biggr]\\
 &\quad=
 (-1)^{d_3}
 \mathcal M_{k+2}(R_{h_3},R_{h_2})
 \mathop{\times}_{R_{h_2}}
 \mathcal M_{\ell+2}(R_{h_2},R_{h_1})\\
 &\hspace{6em}{}
 \mathop{\times}_{R_{h_1}\times L\times\cdots\times L}
 \left(
 S\times\prod_{i=1}^{\ell}P_i
 \times\prod_{j=1}^{k}Q_j
 \right),
\end{aligned}
\]
其中
\[
 d_3=kn(\mu_2-\mu_1+\ell-1).
\]
另一方面，
\[
\begin{aligned}
 &\mathcal M_1
 (R_{h_3},R_{h_1};
 S,P_1,\ldots,P_\ell,Q_1,\ldots,Q_k)\\
 &\quad=
 (-1)^{d_4}
 \mathcal M_{k+\ell+2}(R_{h_3},R_{h_1})
 \mathop{\times}_{R_{h_1}\times L\times\cdots\times L}
 \left(
 S\times\prod_{i=1}^{\ell}P_i
 \times\prod_{j=1}^{k}Q_j
 \right),
\end{aligned}
\]
其中
\[
\begin{aligned}
 d_4={}&
 (k+\ell)(n+1)\mu_1+\mu_3
 +(k+\ell)(n+1)\deg S\\
 &{}+(n+1)
 \sum_{j=1}^{\ell-1}\sum_{i=1}^{j}\deg P_i
 +k(n+1)\sum_{i=1}^{\ell}\deg P_i\\
 &{}+(n+1)
 \sum_{j=1}^{k-1}\sum_{i=1}^{j}\deg Q_i.
\end{aligned}
\]
汇总这些贡献，得到
\[
 \epsilon_3=d_0+d_1+d_2+d_3+d_4
 =(\mu_3+r_3)+1.
\]

最后的断言 (4) 可用同命题~\ref{prop:8-5-1}(1) 类似的方法证明。
\end{proof}

为了定义滤过 $A_\infty$ 双模结构，现在考虑标记点服从标准次序的情形。
\translatornote{源书在下述标准次序中把下边界输入个数写成未在此处定义的
$k_2$；本段随后的定义及全章双模记号均以 $k_0$ 表示 $L^{(0)}$ 一侧输入数，
故统一恢复为 $k_0$。}
也就是说，
\[
 z_0=+\infty,
\]
$z_1,\ldots,z_{k_1}$ 位于 $\mathbb R\times\{1\}$ 上，并满足
\[
 \operatorname{Re}z_1>\cdots>\operatorname{Re}z_{k_1};
\]
又令 $z_{k_1+1}=-\infty$，而
$z_{k_1+2},\ldots,z_{k_1+k_0+1}$ 位于
$\mathbb R\times\{0\}$ 上，并满足
\[
 \operatorname{Re}z_{k_1+2}<\cdots<
 \operatorname{Re}z_{k_1+k_0+1}.
\]
见注~\ref{rem:3-7-23}(2)。计入注~\ref{rem:8-8-9} 与定理
\ref{thm:8-8-10} 后，作如下约定。

\begin{definition}{8.8.11}\label{def:8-8-11}
\begin{enumerate}[label=(\arabic*),leftmargin=2.8em]
\item 设
\[
 P_1^{(1)},\ldots,P_{k_1}^{(1)}
\]
是 $L^{(1)}$ 中的光滑奇异单形，而
\[
 P_1^{(0)},\ldots,P_{k_0}^{(0)}
\]
是 $L^{(0)}$ 中的光滑奇异单形。如下定义第
\ref{subsec:3-7-5} 小节所用模空间的定向：
\[
\begin{aligned}
 &\mathcal M_{k_1,k_0,(+\infty)}
 \bigl(
 L^{(1)},L^{(0)};[h,w],[h',w'];
 \vec P^{\,(1)},S,\vec P^{\,(0)}
 \bigr)\\
 &\quad=
 (-1)^\delta
 \mathcal M_1(R_{h'},R_h;S,P_1,\ldots,P_k),
\end{aligned}
\]
其中
\[
 \delta=(\deg S+1)
 \sum_{i=1}^{k_1}(\deg P_i^{(1)}+1),
\]
并且
\[
 (P_1,\ldots,P_k)
 =
 (P_1^{(1)},\ldots,P_{k_1}^{(1)},
  P_1^{(0)},\ldots,P_{k_0}^{(0)}).
\]
右端模空间用于定义~\ref{def:3-7-71} 中的
$\mathfrak n_{k_1,k_0}$。

\item 对 $\mathfrak n_{0,0}$ 的经典贡献
$\overline{\mathfrak n}_{0,0}$，定义
\[
 \overline{\mathfrak n}_{0,0}
 =
 (-1)^{\dim R_h+\mu(R_h)}\partial
\]
作为 $R_h$ 上链的算子；其中 $\partial$ 是 $R_h$ 上链的通常边界算子。
\end{enumerate}
\end{definition}

在这个新约定（定义~\ref{def:8-8-11}）下，定理
\ref{thm:8-8-10} 蕴含下面的公式
\eqref{eq:8-8-12-1}--\eqref{eq:8-8-12-6}。

\begin{equation}
\begin{aligned}
 &(-1)^{\delta_1}
 \mathcal M_{k_1,k_0,(+\infty)}
 \bigl(
 R_{h_2},R_{h_1};
 P_1^{(1)},\ldots,P_{k_1}^{(1)},
 \partial S,
 P_1^{(0)},\ldots,P_{k_0}^{(0)}
 \bigr)\\
 &\quad\subset
 \partial\mathcal M_{k_1,k_0,(+\infty)}
 \bigl(
 R_{h_2},R_{h_1};
 P_1^{(1)},\ldots,P_{k_1}^{(1)},
 S,
 P_1^{(0)},\ldots,P_{k_0}^{(0)}
 \bigr),
\end{aligned}
\tag{8.8.12.1}\label{eq:8-8-12-1}
\end{equation}
其中
\[
 \delta_1=(\mu_2+r_2)+(\mu_1+r_1)
 +\sum_{j=1}^{k_1}(\deg P_j^{(1)}+1)+1.
\]

\begin{equation}
\begin{aligned}
 &(-1)^{\delta_2}
 \mathcal M_{k_1,k_0,(+\infty)}
 \bigl(
 R_{h_2},R_{h_1};
 P_1^{(1)},\ldots,\partial P_i^{(1)},\ldots,P_{k_1}^{(1)},
 S,P_1^{(0)},\ldots,P_{k_0}^{(0)}
 \bigr)\\
 &\qquad\subset
 \partial\mathcal M_{k_1,k_0,(+\infty)}
 \bigl(
 R_{h_2},R_{h_1};
 P_1^{(1)},\ldots,P_{k_1}^{(1)},
 S,P_1^{(0)},\ldots,P_{k_0}^{(0)}
 \bigr),
\end{aligned}
\tag{8.8.12.2}\label{eq:8-8-12-2}
\end{equation}
其中
\[
 \delta_2=(\mu_2+r_2)+n
 +\sum_{j=1}^{i-1}(\deg P_j^{(1)}+1)+1.
\]

\begin{equation}
\begin{aligned}
 &(-1)^{\delta_3}
 \mathcal M_{k_1,k_0,(+\infty)}
 \bigl(
 R_{h_2},R_{h_1};
 P_1^{(1)},\ldots,P_{k_1}^{(1)},
 S,P_1^{(0)},\ldots,\partial P_i^{(0)},\ldots,P_{k_0}^{(0)}
 \bigr)\\
 &\qquad\subset
 \partial\mathcal M_{k_1,k_0,(+\infty)}
 \bigl(
 R_{h_2},R_{h_1};
 P_1^{(1)},\ldots,P_{k_1}^{(1)},
 S,P_1^{(0)},\ldots,P_{k_0}^{(0)}
 \bigr),
\end{aligned}
\tag{8.8.12.3}\label{eq:8-8-12-3}
\end{equation}
其中\translatornote{源书在 \eqref{eq:8-8-12-3} 的被取边界因子上使用下标
$j$，但随后的符号指数以 $i-1$ 为上限；依
\eqref{eq:8-7-6} 及同组公式的统一记号，把被取边界因子恢复为
$P_i^{(0)}$。}
\[
\begin{aligned}
 \delta_3={}&
 (\mu_2+r_2)+n
 +\sum_{j=1}^{k_1}\deg P_j^{(1)}
 +(\deg S+1)\\
 &{}+\sum_{j=1}^{i-1}(\deg P_j^{(0)}+1)+1.
\end{aligned}
\]

\begin{equation}
\begin{aligned}
 &(-1)^{\delta_4}
 \mathcal M_{i-1,k_0-\ell,(+\infty)}\\[-2pt]
 &\quad
 \Bigl(
 R_{h_3},R_{h_2};
 P_1^{(1)},\ldots,P_{i-1}^{(1)},\\
 &\qquad
 \mathcal M_{k_1-i+1,\ell,(+\infty)}
 \bigl(
 R_{h_2},R_{h_1};
 P_i^{(1)},\ldots,P_{k_1}^{(1)},
 S,P_1^{(0)},\ldots,P_\ell^{(0)}
 \bigr),
 \\[-2pt]
 &\qquad
 P_{\ell+1}^{(0)},\ldots,P_{k_0}^{(0)}
 \Bigr)\\
 &\quad\subset
 \partial\mathcal M_{k_1,k_0,(+\infty)}
 \bigl(
 R_{h_3},R_{h_1};
 P_1^{(1)},\ldots,P_{k_1}^{(1)},
 S,P_1^{(0)},\ldots,P_{k_0}^{(0)}
 \bigr),
\end{aligned}
\tag{8.8.12.4}\label{eq:8-8-12-4}
\end{equation}
其中
\[
 \delta_4=(\mu_3+r_3)
 +\sum_{j=1}^{i-1}(\deg P_j^{(1)}+1)+1.
\]

\begin{equation}
\begin{aligned}
 &(-1)^{\delta_5}
 \mathcal M_{k_1-\ell+i,k_0,(+\infty)}
 \Bigl(
 R_{h_2},R_{h_1};
 P_1^{(1)},\ldots,P_{i-1}^{(1)},\\
 &\qquad
 \mathcal M_{\ell-i+2}
 (\beta;P_i^{(1)},\ldots,P_\ell^{(1)}),
 P_{\ell+1}^{(1)},\ldots,P_{k_1}^{(1)},
 S,P_1^{(0)},\ldots,P_{k_0}^{(0)}
 \Bigr)\\
 &\qquad\subset
 \partial\mathcal M_{k_1,k_0,(+\infty)}
 \bigl(
 R_{h_2},R_{h_1};
 P_1^{(1)},\ldots,P_{k_1}^{(1)},
 S,P_1^{(0)},\ldots,P_{k_0}^{(0)}
 \bigr),
\end{aligned}
\tag{8.8.12.5}\label{eq:8-8-12-5}
\end{equation}
其中
\[
 \delta_5=(\mu_2+r_2)
 +\sum_{j=1}^{i-1}(\deg P_j^{(1)}+1)+1.
\]

\begin{equation}
\begin{aligned}
 &(-1)^{\delta_6}
 \mathcal M_{k_1,k_0-\ell+i,(+\infty)}
 \Bigl(
 R_{h_2},R_{h_1};
 P_1^{(1)},\ldots,P_{k_1}^{(1)},
 S,P_1^{(0)},\ldots,P_{i-1}^{(0)},\\
 &\qquad
 \mathcal M_{\ell-i+2}
 (\beta;P_i^{(0)},\ldots,P_\ell^{(0)}),
 P_{\ell+1}^{(0)},\ldots,P_{k_0}^{(0)}
 \Bigr)\\
 &\qquad\subset
 \partial\mathcal M_{k_1,k_0,(+\infty)}
 \bigl(
 R_{h_2},R_{h_1};
 P_1^{(1)},\ldots,P_{k_1}^{(1)},
 S,P_1^{(0)},\ldots,P_{k_0}^{(0)}
 \bigr),
\end{aligned}
\tag{8.8.12.6}\label{eq:8-8-12-6}
\end{equation}
其中
\[
\begin{aligned}
 \delta_6={}&
 (\mu_2+r_2)
 +\sum_{j=1}^{k_1}(\deg P_j^{(1)}+1)
 +(\deg S+1)\\
 &{}+\sum_{j=1}^{i-1}(\deg P_j^{(0)}+1)+1.
\end{aligned}
\]

因此，$\mathfrak n_{k,\ell}$ 连同符号满足式~\eqref{eq:3-7-2}。
